\documentclass[11pt]{article}
\usepackage{fontspec,amsmath,amssymb,amsthm,hyperref}
\usepackage[margin=1in]{geometry}
\setlength{\emergencystretch}{3em}
\title{Stitched exact segments are the actual global kinetic solution}
\author{LangevinLean mathematical exposition}
\date{}
\begin{document}\maketitle
\noindent Source module: \texttt{KineticLangevin/SegmentGlobal.lean}. Compile this proof document with XeLaTeX or LuaLaTeX. The accompanying HTML pairs the mathematical argument with the exact Lean statements and records the scope of the formalization.


A reference defined one interval at a time must be identified with the actual continuous kinetic process whose stationarity is proved separately. The bridge here uses the true shifted-input equation and uniqueness. It does not assume a Markov transition formula or independence of arbitrarily chosen reference segments.


\section{Identify a selected local solution with a shifted global path}


Start with an actual continuous global solution \(Y\), and suppose the local input is the actual increment \(W_t=B_{a+t}-B_a\). The previously proved time-shift lemma shows \(Y_{a+t}\) satisfies the same kinetic integral equation on the local horizon. The selected local solution starts at \(Y_a\). Lean checks the two initial conditions, the shared additive input, and continuity, then invokes the proved uniqueness theorem for the globally Lipschitz kinetic drift. Their equality follows at every time in the local interval.



\noindent Lean declarations: \texttt{kineticSegmentFlow\_eq\_shift}.

\section{Identify every stitched grid point with the global solution}


The reference chain starts at \(Y_0\). Inductively assume its state at step \(n\) equals \(Y_{nh}\). The next selected segment begins at that same point, so the local-global uniqueness result identifies its endpoint with \(Y_{nh+h}\). Lean performs induction over the actual chain recursion, checks the finite-horizon inequalities using nonnegative step size, and simplifies the real cast of \(n+1\). The result is the exact identity \(R_n=Y_{nh}\), including zero step size.



\noindent Lean declarations: \texttt{kineticReferenceChain\_eq\_global}.

\section{Identify all intermediate times within a reference segment}


Once the starting grid point is identified, the same uniqueness result proves \(Z_n(t)=Y_{nh+t}\) throughout the whole segment. This lets a stationary global phase law supply both the stationary reference position at grid points and the Gaussian reference velocity at intermediate times required by the KLMC local-error proof. No stationary law is assumed for a newly invented process: the stitched reference is proved to be the actual global solution.



\noindent Lean declarations: \texttt{kineticReferenceSegment\_eq\_global}.

\end{document}
